% block: shifted-complex-symmetric-heading page:183 type:section
\subsection{{Shifted Complex Symmetric Linear Systems}}\label{sec:shifted-complex-symmetric-systems}
{Throughout this section, the coefficient matrix $A^{(k)}$ in \eqref{eq:shifted-linear-systems} is assumed to be complex symmetric, i.e., $A^{(k)}=(A^{(k)})^{\top}\ne(A^{(k)})^{\mathrm H}$.}
% block: shifted-cocg-heading page:183 type:subsection
\subsubsection{{The Shifted COCG Method}}\label{subsec:shifted-cocg}
{The derivation of the shifted COCG method \BookCite{184} is the same as in \eqref{eq:shifted-cg-residual-collinearity}--\eqref{eq:shifted-cg-shared-direction}. The only differences between the shifted CG method and the shifted COCG methods are $\alpha_n$ and $\beta_n$. The algorithm of the shifted COCG method is described in Algorithm\ \ref{alg:shifted-cocg}.}
% block: shifted-cocg-algorithm page:183 type:ordinary-paragraph
\begin{algorithm}[H]\caption{{The shifted COCG method (note: $(\bar{\bm a},\bm b)=\bm a^{\top}\bm b$)}}\label{alg:shifted-cocg}\begin{algorithmic}[1]
\Statex\textbf{Input:} {Choose a seed system $s\in S=\{1,2,\ldots,m\}$ and set $\bm r_0^{(s)}=\bm b,\ \beta_{-1}^{(s)}=0$}
\Statex\textbf{Input:} $\bm x_0^{(k)}=\bm p_{-1}^{(k)}=\bm0,\ \pi_0^{(s,k)}=\pi_{-1}^{(s,k)}=1$ {for} $k=1,2,\ldots,m$
\Statex\textbf{Output:} $\bm x_n^{(k)}$ {for} $k=1,2,\ldots,m$
\For{$n=0,1,\ldots$, {until convergence}}
\State $\bm p_n^{(s)}=\bm r_n^{(s)}+\beta_{n-1}^{(s)}\bm p_{n-1}^{(s)}$
\State $\alpha_n^{(s)}=\dfrac{(\bar{\bm r}_n^{(s)},\bm r_n^{(s)})}{(\bar{\bm p}_n^{(s)},A^{(s)}\bm p_n^{(s)})}$
\State $\bm x_{n+1}^{(s)}=\bm x_n^{(s)}+\alpha_n^{(s)}\bm p_n^{(s)}$
\State \{{Begin shifted system}\}
\State {Run lines 6--15 of the shifted CG method (Algorithm\ \ref{alg:shifted-cg}).}
\State \{{End shifted system}\}
\State $\bm r_{n+1}^{(s)}=\bm r_n^{(s)}-\alpha_n^{(s)}A^{(s)}\bm p_n^{(s)}$
\State $\beta_n^{(s)}=\dfrac{(\bar{\bm r}_{n+1}^{(s)},\bm r_{n+1}^{(s)})}{(\bar{\bm r}_n^{(s)},\bm r_n^{(s)})}$
\EndFor\end{algorithmic}\end{algorithm}
% block: shifted-cocg-properties pages:183,184 type:ordinary-paragraph
{Algorithm\ \ref{alg:shifted-cocg} has the following properties:}
\begin{enumerate}
\item {the multiplication of $A^{(s)}$ and a vector is required, but the multiplications of $A^{(k)}$ for the other $k$ and a vector are not required;}
\item {the approximate solution $\bm x_n^{(k)}$ of the shifted COCG method is the same as the $n$th approximate solution of the COCG method (with $\bm x_0=\bm0$) applied to $A^{(k)}\bm x^{(k)}=\bm b$.}
\end{enumerate}
{When matrix $A^{(k)}$ is real symmetric for all $k$, the shifted COCG method is equivalent to the shifted CG method.}
% block: generalized-shifted-cocg page:184 type:ordinary-paragraph
{In \BookCite{190}, the shifted COCG method is generalized to solving the following linear systems:}
\begin{equation}\label{eq:generalized-shifted-systems}(A+\sigma_kB)\bm x^{(k)}=\bm b,\end{equation}
{which are referred to as generalized shifted linear systems. Here $A$ and $B$ are real symmetric and $\sigma_k$'s are complex numbers, and thus the coefficient matrix $A+\sigma_kB$ is complex symmetric. It is easy to see that if $B$ is the identity matrix, then the generalized shifted linear systems \eqref{eq:generalized-shifted-systems} reduce to the shifted linear systems.}
{We note that the shifted COCG method and the other shifted Krylov subspace methods cannot be applied to \eqref{eq:generalized-shifted-systems} because the shift invariance properties do not hold any longer, i.e., $\mathcal K_n(A+\sigma_iB,\bm b)\ne\mathcal K_n(A+\sigma_jB,\bm b)$ for $i\ne j$. One remedy is to consider the following shifted linear systems:}
\[(B^{-1}A+\sigma_k I)\bm x^{(k)}=B^{-1}\bm b,\]
{where $B$ is assumed to be nonsingular. However, $B^{-1}A+\sigma_k I$ is non-Hermitian, and thus one may use shifted Krylov subspace methods for non-Hermitian linear systems. On the other hand, without using such general solvers, it is shown that an algorithm similar to the shifted COCG method is constructed, which is referred to as the generalized shifted COCG method \BookCite{190}. If $B=I$, the generalized shifted COCG method reduces to the shifted COCG method.}
% block: generalized-shifted-applications page:184 type:ordinary-paragraph
{As pointed out in Sect.\ \ref{subsec:electronic-structure}, if the Hamiltonian matrix is real symmetric, then the coefficient matrices $(\epsilon+i\delta)I-H$ of the shifted linear systems are complex symmetric. Thus the shifted COCG method is a method of choice, and applications to computational physics are found in, e.g., \BookCite{73}, \BookCite{110}, \BookCite{111}. An open-source library of the shifted COCG method and related solvers is developed in \BookCite{100}.}
{The generalized shifted linear systems also arise in the following linear systems:}
\begin{equation}\label{eq:quadratic-parameter-linear-system}(\sigma^2A+\sigma B+C)\bm x=\bm b,\end{equation}
{where $A,B,C$ are complex symmetric. Parameterized linear systems of this type are considered in \BookCite{158}, \BookCite{159}, together with their application to structural dynamics. The interesting observation in \BookCite{158}, \S3 is that \eqref{eq:quadratic-parameter-linear-system} can be transformed into the following linear systems:}
\begin{equation}\label{eq:quadratic-system-shifted-linearization}\left(\begin{bmatrix}B&C\\C^{\top}&O\end{bmatrix}+\sigma\begin{bmatrix}A&O\\O&-C^{\top}\end{bmatrix}\right)\begin{bmatrix}\bm y\\\bm x\end{bmatrix}=\begin{bmatrix}\bm b\\\bm0\end{bmatrix}\end{equation}
{if $C$ is nonsingular. Since $A,B,C$ are complex symmetric, the coefficient matrices of \eqref{eq:quadratic-system-shifted-linearization} are complex symmetric generalized shifted linear systems. Thus, the generalized shifted COCG method can also be used to solve \eqref{eq:quadratic-parameter-linear-system}.}
% block: seed-switching-heading page:185 type:subsection
\paragraph{{Seed Switching Technique}}\label{subsubsec:seed-switching}\mbox{}\par
{We can see from Algorithm\ \ref{alg:shifted-cocg} (see also line 8 of Algorithm\ \ref{alg:shifted-cg}) that if $|\pi_n^{(s,k)}|=|R_n^{(s)}(\sigma_s-\sigma_k)|\ge1$, then $\|\bm r_n^{(k)}\|\le\|\bm r_n^{(s)}\|$. Hence, if we could find a seed system such that $|R_n^{(s)}(\sigma_s-\sigma_k)|\ge1$, then all the shifted systems could be solved. However, it is extremely difficult to find such a system in advance except for some special cases discussed in \BookCite{67}, Corollary 1. The seed switching technique \BookCite{166} will avoid such a problem, and consists of the following steps:}
\begin{enumerate}[label=(S\arabic*)]
\item {Choose a seed system, and then start Algorithm\ \ref{alg:shifted-cocg}.}
\item {If the seed system was solved at the $n$th iteration, then find a new one that remained unsolved.}
\item {Start Algorithm\ \ref{alg:shifted-cocg} from the $(n+1)$th iteration using the new seed system.}
\end{enumerate}
% block: seed-switching-first-stage page:185 type:ordinary-paragraph
{In (S2), one of the criteria for choosing the new seed system $\widetilde s$ may be}
\[\widetilde s=\operatorname*{arg\,max}_{i\in I}\{\|\bm r_n^{(i)}\|\},\]
{where $I$ denotes an index set of the unsolved linear systems. In (S3), we need two steps to switch the old seed system to the new one. First, compute}
\[\pi_{n+1}^{(s,\widetilde s)}=R_{n+1}^{(s)}(\sigma_s-\sigma_{\widetilde s}),\quad\beta_n^{(\widetilde s)}=\left(\pi_n^{(s,\widetilde s)}/\pi_{n+1}^{(s,\widetilde s)}\right)^2\beta_n\]
{to obtain $\bm r_{n+1}^{(\widetilde s)}$ and $\beta_n^{(\widetilde s)}\bm p_n^{(\widetilde s)}$. Since it follows from $\bm r_{n+1}^{(\widetilde s)}+\beta_n^{(\widetilde s)}\bm p_n^{(\widetilde s)}$ that we obtain $\bm p_{n+1}^{(\widetilde s)}$, one can start the COCG method solving the system $(A+\sigma_{\widetilde s}I)\bm x^{(\widetilde s)}=\bm b$ from the $(n+1)$th iteration step. Second, to solve the remaining linear systems by using the new seed $\widetilde s$, the parameters $\alpha_{n+1}^{(i)}$ and $\beta_n^{(i)}$ must be generated from the new seed. We see that they can readily be generated by the following polynomial:}
\[\pi_{n+1}^{(\widetilde s,i)}=R_{n+1}^{(\widetilde s)}(\sigma_{\widetilde s}-\sigma_i)\quad\text{for all }i\in I.\]
% block: seed-switching-second-stage page:185 type:ordinary-paragraph
{To obtain the above polynomial, one needs to compute}
\[\alpha_i^{(\widetilde s)}=\left(\pi_i^{(s,\widetilde s)}/\pi_{i+1}^{(s,\widetilde s)}\right)\alpha_i,\quad\beta_j^{(\widetilde s)}=\left(\pi_j^{(s,\widetilde s)}/\pi_{j+1}^{(s,\widetilde s)}\right)^2\beta_j\]
{for $i=0,\ldots,n$, $j=0,\ldots,n-1$. Hence, the switching strategy requires only scalar operations, and moreover we can see that if breakdown does not occur, iterating the process from (S2) to (S3) enables us to keep solving the systems without losing the dimension of the Krylov subspace that has been generated until the last switching. The seed switching technique enables us to be free from the problem of the choice of the seed system, and can also be applied to the shifted CG method in Sect.\ \ref{subsec:shifted-cg} and the shifted BiCG method in Sect.\ \ref{subsec:shifted-bicg}, since the residual vectors of the shifted CG (BiCG) method for shifted systems are true CG (BiCG) residuals.}
% block: shifted-cocr-heading page:186 type:subsection
\subsubsection{{The Shifted COCR Method}}\label{subsec:shifted-cocr}
{The derivation of the shifted COCR method \BookCite{172} is also the same as in \eqref{eq:shifted-cg-residual-collinearity}--\eqref{eq:shifted-cg-shared-direction}. The only differences between the shifted CR method and the shifted COCR methods are $\alpha_n$ and $\beta_n$. The algorithm of the shifted COCR method is described in Algorithm\ \ref{alg:shifted-cocr}.}
% block: shifted-cocr-properties page:186 type:ordinary-paragraph
{Algorithm\ \ref{alg:shifted-cocr} has the following properties:}
\begin{enumerate}
\item {the multiplication of $A^{(s)}$ and a vector is required, but the multiplications of $A^{(k)}$ for the other $k$ and a vector are not required, which is the same property as that of the shifted COCG method;}
\item {unlike the shifted COCG method, $\bm x_n^{(k)}$ of the shifted COCR method is not the same as the $n$th approximate solution of the COCR method (with $\bm x_0=\bm0$) applied to $A^{(k)}\bm x^{(k)}=\bm b$.}
\end{enumerate}
% block: shifted-cocr-algorithm page:186 type:ordinary-paragraph
\begin{algorithm}[H]\caption{{The shifted COCR method (note: $(\bar{\bm a},\bm b)=\bm a^{\top}\bm b$)}}\label{alg:shifted-cocr}\begin{algorithmic}[1]
\Statex\textbf{Input:} {Choose a seed system $s\in S=\{1,2,\ldots,m\}$ and set $\bm r_0^{(s)}=\bm b,\ \beta_{-1}^{(s)}=0$}
\Statex\textbf{Input:} $\bm x_0^{(k)}=\bm p_{-1}^{(k)}=\bm0,\ \pi_0^{(s,k)}=\pi_{-1}^{(s,k)}=1$ {for} $k=1,2,\ldots,m$
\Statex\textbf{Output:} $\bm x_n^{(k)}$ {for} $k=1,2,\ldots,m$
\For{$n=0,1,\ldots$, {until convergence}}
\State $\bm p_n^{(s)}=\bm r_n^{(s)}+\beta_{n-1}^{(s)}\bm p_{n-1}^{(s)}$
\State $(A^{(s)}\bm p_n^{(s)}=A^{(s)}\bm r_n^{(s)}+\beta_{n-1}^{(s)}A^{(s)}\bm p_{n-1}^{(s)})$
\State $\alpha_n^{(s)}=\dfrac{(\bar{\bm r}_n^{(s)},A^{(s)}\bm r_n^{(s)})}{(\overline{A^{(\sigma_s)}\bm p_n^{(s)}},A^{(s)}\bm p_n^{(s)})}$
\State $\bm x_{n+1}^{(s)}=\bm x_n^{(s)}+\alpha_n^{(s)}\bm p_n^{(s)}$
\State \{{Begin shifted system}\}
\State {Run lines 6--15 of the shifted CG method (Algorithm\ \ref{alg:shifted-cg}).}
\State \{{End shifted system}\}
\State $\bm r_{n+1}^{(s)}=\bm r_n^{(s)}-\alpha_n^{(s)}A^{(s)}\bm p_n^{(s)}$
\State $\beta_n^{(s)}=\dfrac{(\bar{\bm r}_{n+1}^{(s)},A^{(s)}\bm r_{n+1}^{(s)})}{(\bar{\bm r}_n^{(s)},A^{(s)}\bm r_n^{(s)})}$
\EndFor\end{algorithmic}\end{algorithm}
% block: shifted-qmr-sym-heading page:186 type:subsection
{When matrix $A^{(k)}$ is real symmetric for all $k$, the shifted COCR method is equivalent to the shifted CR method.}
\subsubsection{{The Shifted QMR\_SYM Method}}\label{subsec:shifted-qmr-sym}
{The shifted QMR method for solving non-Hermitian linear systems was proposed in \BookCite{62}. As is known in \BookCite{61}, the shifted QMR\_SYM method is a simplification of the shifted QMR method, and the derivation of the shifted QMR\_SYM is similar to the shifted MINRES method.}
% block: shifted-qmr-sym-projection page:187 type:ordinary-paragraph
{From the complex symmetric Lanczos process in Sect.\ \ref{subsec:complex-symmetric-lanczos}, $\bm x_n\in\mathcal K_n(A,\bm b)$ can be rewritten as}
\begin{equation}\label{eq:shifted-qmr-sym-coordinate-solution}\bm x_n^{(k)}=V_n\bm y_n^{(k)},\quad k=1,2,\ldots,m,\end{equation}
{where $\bm y_n^{(k)}\in\mathbb C^n$. The corresponding residual vectors for \eqref{eq:shifted-linear-systems} are $\bm r_n^{(k)}:=\bm b-(A+\sigma_k I)\bm x_n^{(k)}$. From \eqref{eq:shifted-qmr-sym-coordinate-solution} and the matrix form of the complex symmetric Lanczos process \eqref{eq:complex-lanczos-matrix-relation}, it follows that}
\begin{equation}\label{eq:shifted-qmr-sym-projected-residual}\begin{aligned}\bm r_n^{(k)}&=\bm b-(A+\sigma_k I)V_n\bm y_n^{(k)}\\&=V_{n+1}g_1\bm e_1-(AV_n+\sigma_kV_n)\bm y_n^{(k)}\\&=V_{n+1}(g_1\bm e_1-T_{n+1,n}^{(k)}\bm y_n^{(k)}),\quad T_{n+1,n}^{(k)}:=T_{n+1,n}+\sigma_k\begin{bmatrix}I_n\\\bm0^{\top}\end{bmatrix}.\end{aligned}\end{equation}
{Here, $\bm e_1=(1,0,\ldots,0)^{\top}$ is the first unit vector and $g_1=(\bar{\bm b},\bm b)^{1/2}$. As well as the QMR\_SYM method in Sect.\ \ref{subsec:qmr-sym}, $\bm y_n^{(k)}$ is determined by solving the following minimization problem.}
\begin{equation}\label{eq:shifted-qmr-sym-least-squares}\min_{\bm y_n^{(k)}\in\mathbb C^n}\|g_1\bm e_1-T_{n+1,n}^{(k)}\bm y_n^{(k)}\|.\end{equation}
{The minimization problem \eqref{eq:shifted-qmr-sym-least-squares} can be solved by Givens rotations that are described in Sect.\ \ref{subsec:minres}. The algorithm of the shifted QMR\_SYM method is shown in Algorithm\ \ref{alg:shifted-qmr-sym}.}
% block: shifted-qmr-sym-properties page:187 type:ordinary-paragraph
{Algorithm\ \ref{alg:shifted-qmr-sym} holds the following properties:}
\begin{enumerate}
\item {the multiplication of $A^{(s)}$ and a vector is required, but the multiplications of $A^{(k)}$ for the other $k$ and a vector are not required, which is the same property as those of the shifted CG method and the shifted CR method;}
\item {the approximate solution $\bm x_n^{(k)}$ of the shifted QMR\_SYM method is the same as the $n$th approximate solution of the QMR\_SYM method (with $\bm x_0=\bm0$) applied to $A^{(k)}\bm x^{(k)}=\bm b$.}
\end{enumerate}
% block: shifted-qmr-sym-algorithm page:187 type:ordinary-paragraph
\begin{algorithm}[H]\caption{{The shifted QMR\_SYM method (note: $(\bar{\bm a},\bm b)=\bm a^{\top}\bm b$)}}\label{alg:shifted-qmr-sym}\begin{algorithmic}[1]
\Statex\textbf{Input:} $\bm x_0^{(k)}=\bm p_{-1}^{(k)}=\bm p_0^{(k)}=\bm0,\ \bm v_1=\bm b/(\bar{\bm b},\bm b)^{1/2},\ g_1^{(k)}=(\bar{\bm b},\bm b)^{1/2},\ \beta_{-1}=0$
\Statex\textbf{Output:} $\bm x_n^{(k)}$ {for} $l=1,2,\ldots,m$
\For{$n=1,2,\ldots$}
\State ({The complex symmetric Lanczos process})
\State $\alpha_n=(\bar{\bm v}_n,A\bm v_n)$
\State $\widetilde{\bm v}_{n+1}=A\bm v_n-\alpha_n\bm v_n-\beta_{n-1}\bm v_{n-1}$
\State $\beta_n=(\overline{\widetilde{\bm v}_{n+1}},\widetilde{\bm v}_{n+1})^{1/2}$
\State $\bm v_{n+1}=\widetilde{\bm v}_{n+1}/\beta_n$
\State $t_{n-1,n}^{(\ell)}=\beta_{n-1},\ t_{n,n}^{(\ell)}=\alpha_n+\sigma_{\ell},\ t_{n+1,n}^{(\ell)}=\beta_n$
\State ({Solve least-squares problems by Givens rotations})
\State {Run lines 9--24 of the shifted MINRES method (Algorithm\ \ref{alg:shifted-minres}).}
\EndFor\end{algorithmic}\end{algorithm}
% block: hermitian-complex-shift-qmr page:188 type:ordinary-paragraph
{In \BookCite{60}, the following shifted linear systems are considered:}
\begin{equation}\label{eq:hermitian-matrix-complex-shift}(A+\sigma_k I)\bm x^{(k)}=\bm b,\end{equation}
{where $A$ is Hermitian and $\sigma_k$'s are complex numbers. Thus the coefficient matrix $(A+\sigma_k I)$ is neither Hermitian nor complex symmetric. In this case, one may use the shifted QMR method for non-Hermitian shifted linear systems. On the other hand, there is an algorithm in \BookCite{60} that is more efficient than the shifted QMR method so that the multiplication of $(A+\sigma_k I)^{\mathrm H}$ and a vector is not required, i.e., the number of required matrix--vector multiplications is only one per each iteration step. The key idea is to generate a Krylov subspace of $\mathcal K_n(A,\bm b)$ by the Lanczos process, and the basis vectors are used to solve \eqref{eq:hermitian-matrix-complex-shift}. Note that if we generate $\mathcal K_n(A+\sigma I,\bm b)$ with $\sigma$ being a complex number instead of $\mathcal K_n(A,\bm b)$, then we need the bi-Lanczos process whose computational cost is about twice as large as that of the Lanczos process when the matrix--vector multiplication is the most time-consuming part.}
% block: generalized-shifted-qmr page:188 type:ordinary-paragraph
{In \BookCite{168}, the shifted QMR\_SYM method is generalized to solving (complex symmetric) generalized shifted linear systems \eqref{eq:generalized-shifted-systems}. Similar to the generalized shifted COCG method as mentioned in Sect.\ \ref{subsec:shifted-cocg}, the generalized shifted QMR\_SYM method \BookCite{168} reduces to the shifted QMR\_SYM method when $B=I$. Furthermore, based on the shifted weighted QMR\_SYM method in \BookCite{167} for complex symmetric shifted linear systems, the corresponding algorithm for solving complex symmetric generalized shifted linear systems is proposed in \BookCite{168}.}
